\documentclass[12pt]{article}
\usepackage[T1]{fontenc}
\usepackage[utf8]{inputenc}
\usepackage{lmodern}
\usepackage{textcomp}
\usepackage{enumitem}
\usepackage{geometry}
\geometry{margin=1in}
\begin{document}
\begin{center}
Answer Key - Philosophy - Undergraduate Level Assessment 21 \
\small (Answers are ordered exactly as in the question paper.)
\end{center}

\section*{Multiple Choice}
\begin{enumerate}[label=\arabic*.]
\item \textbf{Answer:} C
\\ \textit{Options:} A) Orthodox; B) Sephardic; C) Conservative; D) Hasidic; E) Reconstructionism
\\ \textit{Note:} The Conservative branch of Judaism, founded by Zacharias Frankel, is known for its "Positive-Historical Judaism" because it emphasizes a historical understanding of Jewish law and tradition while allowing for adaptation and change in response to contemporary circumstances.
\item \textbf{Answer:} A
\\ \textit{Options:} A) suggest that aid agencies are the only solution to global poverty.; B) abandon the world's poor to their unjust predicament.; C) argue that aid agencies are inherently unjust.; D) increase the amount of aid given to the world's poor.; E) ignore the issue of global economic justice.
\\ \textit{Note:} Ashford argues that while concerns about dependency on aid agencies are valid, they do not justify the notion that these agencies are the sole solution to global poverty, as addressing poverty requires a multifaceted approach that includes empowering individuals and communities rather than solely relying on external aid.
\item \textbf{Answer:} A
\\ \textit{Options:} A) the artist's nationality; B) its power to evoke strong emotions; C) social meaning; D) the use of geometric shapes; E) the size of the artwork
\\ \textit{Note:} The artist's nationality is an important formal characteristic of art because it often influences the themes, styles, and techniques found in their work, reflecting cultural identity and social context.
\item \textbf{Answer:} D
\\ \textit{Options:} A) the false dilemma fallacy.; B) the appeal to authority fallacy.; C) the straw man fallacy.; D) the naturalistic fallacy.; E) the ad hominem fallacy.
\\ \textit{Note:} Moore asserts that many past philosophers mistakenly equate "good" with natural properties or empirical facts, committing the naturalistic fallacy by trying to define ethical terms in terms of non-ethical ones.
\item \textbf{Answer:} A
\\ \textit{Options:} A) we believe in the possibility of redemption and change.; B) there are some penalties worse than death.; C) we are unable to administer justice effectively.; D) there is no punishment that is proportional to murder.; E) there is always a possibility of judicial errors.
\\ \textit{Note:} This answer is correct because Nathanson emphasizes that abolishing the death penalty reflects a societal commitment to the belief that individuals can reform and be rehabilitated, rather than being permanently condemned for their past actions.
\end{enumerate}

\section*{True / False}
\begin{enumerate}[label=\arabic*.]
\setcounter{enumi}{5}
\item \textbf{Answer:} False
\\ \textit{Options:} A) True; B) False
\\ \textit{Note:} Gauthier's foundational crisis in morality primarily concerns the rational justification of moral principles and social cooperation, rather than being centered on divine command theory, which emphasizes moral obligations based on divine authority.
\item \textbf{Answer:} False
\\ \textit{Options:} A) True; B) False
\\ \textit{Note:} Velleman argues against euthanasia by emphasizing that it undermines the intrinsic value of life and the moral obligations of care, suggesting that such practices do not enhance, but rather compromise, patient care and autonomy.
\item \textbf{Answer:} True
\\ \textit{Options:} A) True; B) False
\\ \textit{Note:} The expression Gba is true because, in predicate logic, the function G represents the action of greeting, and the variables b and a correspond to the individuals Ben and Alexis, respectively, correctly capturing the intended meaning of "Ben greets Alexis."
\item \textbf{Answer:} True
\\ \textit{Options:} A) True; B) False
\\ \textit{Note:} The objective truth of a claim is determined by its correspondence to reality, independent of the personal beliefs, actions, or inconsistencies of the individual making the claim.
\item \textbf{Answer:} False
\\ \textit{Options:} A) True; B) False
\\ \textit{Note:} Moore's ideal utilitarianism emphasizes maximizing intrinsic goods, such as beauty and friendship, rather than merely focusing on the generation of wealth as the primary criterion for determining the right action.
\end{enumerate}

\section*{Long Form Response}
\begin{enumerate}[label=\arabic*.]
\setcounter{enumi}{10}
\item \textbf{Answer:} The fallacy of appeal to loyalty occurs when an individual is urged to act in a certain way primarily due to their allegiance to a person or group, rather than through rational consideration of the situation. This form of reasoning can significantly influence decision-making by compelling individuals to prioritize loyalty over objective analysis, often leading to biased or harmful outcomes. For instance, a person may support unethical practices within their organization simply to remain loyal to their colleagues, thereby compromising their moral judgment. The ethical implications of this fallacy are profound, as it raises questions about the integrity of personal values and the potential for loyalty to perpetuate injustice. Ultimately, while loyalty is a valuable trait, it should not overshadow the necessity of critical thinking and ethical reasoning in decision-making processes.
\\ \textit{Note:} The answer correctly identifies that the fallacy of appeal to loyalty pressures individuals to make decisions based on allegiance rather than rationality, which can result in unethical behavior and detrimental consequences.
\item \textbf{Answer:} The relationship between rights and duties is fundamentally reciprocal, as the existence of one often implies the necessity of the other. In the case of John's duty to return the car to Mary, John's obligation stems from the right that Mary has to her property. When John borrowed the car, he entered into an implicit agreement that conferred upon him the right to use the vehicle while simultaneously creating a duty to return it. This duty ensures that Mary's rights are respected and upheld, illustrating how rights and duties coexist in a moral framework where the fulfillment of one supports the other. Thus, John's action of returning the car not only fulfills his obligation but also reinforces Mary's ownership rights, highlighting the interconnectedness of rights and duties in social interactions.
\\ \textit{Note:} correct because it illustrates the reciprocal nature of rights and duties, demonstrating that John's duty to return Mary's car directly arises from her right to her property, thereby ensuring the respect and fulfillment of that right.
\end{enumerate}

\vfill
\noindent\textit{End of Answer Key}
\end{document}